\section{Aquinas: Sign, Sacrament, and Sacrifice}

Aquinas gives conceptual precision to what the Fathers confess liturgically.
The Eucharist is a memorial of the Lord's Passion, the fulfillment of figures,
and a pledge of future glory (\emph{Summa theologiae} III, q.~73, a.~4).
These temporal directions coexist in one sacrament without erasing the order
among them.

In treating the rite as sacrifice, Aquinas says that celebration of the
sacrament is an image representing Christ's Passion, while the consecrated
sacrament itself contains Christ who suffered (III, q.~83, a.~1). ``Image''
does not mean unreality, because the sign contains the signified Christ
substantially. ``Contains'' does not mean a second Passion, because Christ now
exists risen and impassible. The whole risen Christ is present under either
species; the separate consecrations sacramentally represent the separation of
Body and Blood in death without dividing Christ or making him suffer now (III,
q.~76, a.~2). The
sacramental mode makes present the victim and represents the Passion; it does
not reproduce the Passion's bloody mode.

This distinction is exactly what chronological language tends to blur. A
person transported to Calvary would encounter the suffering Christ under the
historical appearances proper to that event. At Mass the faithful encounter
the whole risen Christ substantially present under either species,
and his sacrifice in an unbloody sacramental manner. The difference is not a
deficiency to be explained away. It belongs to what the Eucharist is.

\subsection*{Cause, representation, and application}

Aquinas also distinguishes the sacrament's relation to past, present, and
future: it commemorates the Passion, contains present grace and unity, and
prefigures future glory (III, q.~73, a.~4). The Passion is the cause of the
Eucharist's sacrificial efficacy, but causality does not require clock-time
identity. A past historical cause can remain effective; the sacrament is the
divinely instituted instrument by which its fruit is communicated here and
now.

In III, q.~79, a.~1, Aquinas treats the sacrament's principal effect as
spiritual nourishment and union with Christ through grace and charity. In III,
q.~83, a.~1, he treats the celebration as a representation of the Passion and
as an offering for living and dead. These distinctions matter: Christ's
objective presence and sacrifice do not depend on a communicant's imaginative
ability to reconstruct Calvary.

